% Добавление fleqn для выравнивания по левому краю и leqno для номеров слева
% 添加 fleqn 选项使公式左对齐，leqno 选项使编号显示在左侧
\documentclass[a4paper, 12pt, fleqn, leqno]{article}

% Настройка кодировок для русского языка / 俄语字符编码设置
\usepackage[T2A]{fontenc}
\usepackage[utf8]{inputenc}
\usepackage[russian]{babel}

% Подключение пакетов amsmath и bm / 加载数学增强宏包和加粗宏包
\usepackage{amsmath}
\usepackage{bm}

\begin{document}
	
	% Упражнение 1: Сравнение inline и display math / 练习 1：行内与行间模式对比 
	% Сравнивает компактную строчную ($\$$) и полноразмерную выключную ($\text{\[...\]}$) верстку формул.
	\section*{Сравнение режимов}
	Формула в текстовом режиме (inline): $\sum_{i=1}^{n} i^2 = \frac{n(n+1)(2n+1)}{6}$. 
	% 注意行内模式下求和符号的上下限会被压缩
	
	Формула в выключном режиме (display):	
	\[
	\sum_{i=1}^{n} i^2 = \frac{n(n+1)(2n+1)}{6}
	\]
	% 在独立行模式下，符号会以更舒展的样式显示Ы
	
	% Упражнение 2: Греческие буквы / 练习 2：希腊字母
	% Показывает различные начертания (курсивное, прямое, жирное) и их вложенность.
	\section*{Греческие буквы}
	Нижний регистр : $\alpha, \beta, \gamma, \delta, \theta, \pi, \omega$. \\
	Верхний регистр : $\Gamma, \Delta, \Theta, \Pi, \Omega$.
	
	% Упражнение 3: Шрифты и вложенность / 练习 3：数学字体与嵌套
	% Демонстрирует написание греческих букв и выделение символов жирным через bm.
	\section*{Шрифты в математике}
	Обычный курсив : $abcXYZ$ \\
	Жирный шрифт : $\mathbf{abcXYZ}$ \\
	Прямой шрифт : $\mathrm{abcXYZ}$ \\
	Использование \texttt{bm} для символов: $\alpha$ vs $\bm{\alpha}$.
	
	% Эксперимент с вложенностью (嵌套实验)
	Пример вложенности шрифтов: $\mathrm{\mathbf{BoldUprightText}}$
	
	% Упражнение 4: Нумерованные уравнения / 练习 4：带编号的等式
	% Создает нумерованное уравнение с индексом слева согласно опции leqno.
	\section*{Нумерованные уравнения}
	Благодаря опции \texttt{leqno}, номер этой формулы будет слева:
	\begin{equation}
		E = mc^2
	\end{equation}
	
\end{document}